\documentclass[11pt,a4paper]{article}
\usepackage[T1]{fontenc}
\usepackage[polish]{babel}
\usepackage{geometry}
\usepackage{listings}
\usepackage{xcolor}
\usepackage{hyperref}
\usepackage{booktabs}
\usepackage{float}

\geometry{margin=2cm}
\setlength{\parskip}{0.4cm}

\lstset{
    language=Python,
    basicstyle=\ttfamily\small,
    keywordstyle=\color{blue},
    commentstyle=\color{gray},
    stringstyle=\color{red},
    breaklines=true,
    backgroundcolor=\color{lightgray!10},
    frame=single,
    captionpos=b
}

\title{\textbf{WIG20 Agent System} \\ \normalsize Dokumentacja Techniczna}
\author{}
\date{\today}

\begin{document}

\maketitle

\tableofcontents

\newpage

\section{Streszczenie}

Projekt WIG20 Agent to zaawansowany system do automatycznego pobierania, przechowywania i analizy danych gieldowych z indeksu WIG20. Rozwiazanie integruje nowoczesne technologie web scrapingu, relacyjne bazy danych (MySQL) oraz model sztucznej inteligencji (Ollama) do odpowiadania na pytania uzytkownikow w jezyku naturalnym.

System realizuje kompleksowy przepyw danych zaczynajacy sie od pobrania informacji z portalu finansowego Bankier.pl, poprzez przechowywanie w bazie danych, az do prezentacji wynikow uzytkownikowi za pomoca inteligentnego agenta AI.

\textbf{Glowne mozliwosci:}
\begin{itemize}
    \item Automatyczne i okresowe pobieranie cen akcji ze strony Bankier.pl
    \item Przechowywanie danych w strukturyzowanej bazie MySQL z historia aktualizacji
    \item REST API do wysylania zapytan i komunikacji z systemem
    \item Inteligentne odpowiedzi na pytania uzytkownikow za pomoca modelu Ollama
    \item Obsuga wielu jezykow (polski i angielski)
\end{itemize}

\section{Cel i Zakres Projektu}

Celem projektu jest stworzenie narzedzia do:

\begin{enumerate}
    \item \textbf{Automatyzacji pobierania danych} - system periodycznie skanuje strone Bankier.pl i ekstrahuje dane o cenach akcji spółek wchodzacych w sklad indeksu WIG20
    \item \textbf{Przechowywania danych} - wszystkie zebrane dane sa zapisywane w bazie MySQL z metadanymi czasowymi
    \item \textbf{Udostepniania interfejsu} - REST API umozliwia komunikacje z systemem poprzez zapytania HTTP
    \item \textbf{Analizy inteligentnej} - model AI (Ollama) przetwaza pytania uzytkownikow i generuje odpowiedzi oparte na aktualnych danych
\end{enumerate}

\section{Funkcjonalnosc}

System oferuje nastepujace funkcjonalnosci:

\subsection{Pobieranie Danych}

Automatyczne pobranie danych z Bankier.pl obejmuje:
\begin{itemize}
    \item Symbole handlowe spółek (np. PKO, PKNORLEN, ORLEN)
    \item Ceny akcji wyrazone w PLN
    \item Kapitalizacje rynkowe spółek
    \item Znaczniki czasowe (timestamp) kazdego pobrania
\end{itemize}

\subsection{Baza Danych}

MySQL przechowuje:
\begin{itemize}
    \item Tabele \texttt{Symbols} - liste wszystkich spółek z ich kapitalizacja
    \item Tabele \texttt{Prices} - historie cen z datami aktualizacji
    \item Unikatowe indeksy zapobiegajace duplikatom
    \item Relacje miedzy tabelami poprzez klucze obce
\end{itemize}

\subsection{API REST}

Endpoint POST /ask umozliwia:
\begin{itemize}
    \item Wysylanie zapytan w jezyku naturalnym
    \item Klasyfikacje intencji zapytania (porownanie, analiza, przeglad)
    \item Pobieranie odpowiednich danych z bazy
    \item Zwracanie wzbogaconych odpowiedzi
\end{itemize}

\subsection{AI i Przetwarzanie Jezyka}

Ollama z modelem Qwen2.5 odpowiada za:
\begin{itemize}
    \item Rozumienie zapytan w jezyku polskim i angielskim
    \item Porownywanie spółek na podstawie danych
    \item Analize trendow rynkowych
    \item Generowanie naturalnych odpowiedzi dla uzytkownika
\end{itemize}

\section{Architektura Systemu}

\subsection{Komponenty}

System sklada sie z czterech glownych komponentów:

\subsubsection{scraper.py - Pobieranie Danych}

Modul odpowiedzialny za:
\begin{itemize}
    \item Wysylanie zapytan HTTP do Bankier.pl
    \item Parsowanie odpowiedzi HTML przy uzyciu BeautifulSoup
    \item Ekstrakcje danych za pomoca wyrazow regularnych (regex)
    \item Obsluge bledów i ponowne probowanie w przypadku niepowodzenia
    \item Emulacje przegladarki poprzez User-Agent (unikanie blokady)
\end{itemize}

\subsubsection{db\_updater.py - Aktualizacja Bazy}

Warstwa bazodanowa zajmujaca sie:
\begin{itemize}
    \item Tworzeniem tabel (Symbols, Prices) przy pierwszym uruchomieniu
    \item Zapisywaniem danych z scrapera do MySQL
    \item Operacjami UPSERT dla idempotentnosci
    \item Obsuga duplikatow poprzez unikatowe indeksy
    \item Logowaniem i raportowaniem bledów
\end{itemize}

\subsubsection{main.py - REST API}

FastAPI serwer dostarczajacy:
\begin{itemize}
    \item Endpoint \texttt{POST /ask} do wysylania pytań
    \item Klasyfikacje typu zapytania (porownanie/analiza/przeglad)
    \item Logike pobierania danych z bazy na podstawie zapytania
    \item Integracje z modelem Ollama
    \item Dokumentacje API (Swagger UI na \texttt{/docs})
\end{itemize}

\subsubsection{mcp\_server.py - Model Context Protocol}

Zaawansowany serwer komunikacyjny oferujacy:
\begin{itemize}
    \item Zasoby (resources): \texttt{stock://wig20/all}, \texttt{stock://wig20/prices}
    \item Narzedzia (tools): \texttt{set\_price\_alert}
    \item Transport Server-Sent Events (SSE) dla aktualizacji real-time
    \item Wsparcie dla zaawansowanych integracji
\end{itemize}

\section{Technologie}

\subsection{Stack Techniczny}

\begin{itemize}
    \item \textbf{Python 3.11} - jezyk programowania (wydajny, zrodla open-source)
    \item \textbf{FastAPI} - nowoczesny framework do budowy REST API
    \item \textbf{Uvicorn} - szybki serwer ASGI
    \item \textbf{MySQL 8.0} - relacyjna baza danych
    \item \textbf{BeautifulSoup4} - parsing HTML
    \item \textbf{requests} - HTTP client dla Pythona
    \item \textbf{Pydantic} - walidacja danych i type hints
    \item \textbf{Ollama} - framework dla modeli LLM
    \item \textbf{Qwen2.5} - model AI/LLM
    \item \textbf{Docker Compose} - orkiestracja kontenerow
\end{itemize}

\subsection{Wymagania Systemowe}

\begin{itemize}
    \item Python 3.11 lub nowszy
    \item Docker 20.10+
    \item Docker Compose 1.29+
    \item RAM: minimum 4GB (zalecane 8GB dla Ollama)
    \item Dysk: minimum 10GB na baze danych i model AI
    \item Dostep do internetu dla pobrania modelu Qwen2.5
\end{itemize}

\section{Deployment i Infrastruktura}

\subsection{Docker Compose}

System uruchamia sie w czterech izolowanych kontenerach:

\begin{table}[H]
\centering
\begin{tabular}{|l|l|l|l|}
\hline
\textbf{Nazwa} & \textbf{Obraz} & \textbf{Port} & \textbf{Opis} \\
\hline
db & mysql:8.0 & 3306 & Baza danych MySQL \\
ollama & ollama/ollama & 11434 & Model AI/LLM \\
api & local/wig20:latest & 8000 & FastAPI service \\
mcp & local/wig20:latest & - & MCP Server \\
\hline
\end{tabular}
\caption{Serwisy Docker Compose}
\end{table}

\subsection{Health Checks i Zaleznosci}

System implementuje:
\begin{itemize}
    \item Health check dla MySQL: \texttt{mysqladmin ping -h localhost}
    \item Timeout: 20 sekund
    \item Liczba prob: 10 ponownych prob
    \item Zaleznosci: API i MCP czekaja na gotowość bazy danych
    \item Uruchamianie w okreslonym porzadku dla spójnosci
\end{itemize}

\subsection{Konfiguracja Zmiennych Srodowiskowych}

Plik \texttt{.env} zawiera:
\begin{lstlisting}
MYSQL_ROOT_PASSWORD=haslo_root
MYSQL_DATABASE=wig20
MYSQL_USER=agent
MYSQL_PASSWORD=haslo_agenta
\end{lstlisting}

\section{Przepyw Danych i Architektura}

\subsection{Sekwencja Zdarzen}

Pełny przepyw danych:

\begin{enumerate}
    \item \textbf{Pobranie}: Scraper wysyla zapytanie do Bankier.pl
    \item \textbf{Parsing}: BeautifulSoup ekstrahuje HTML i wydobywa dane
    \item \textbf{Zapis}: db\_updater zapisuje dane do MySQL z deduplicacja
    \item \textbf{Zapytanie}: Uzytkownik wysyla POST na /ask
    \item \textbf{Klasyfikacja}: API okresla typ zapytania
    \item \textbf{Pobieranie}: Zapytania SQL pobieraja relewalntne dane
    \item \textbf{AI}: Dane wysylane do Ollamy z promptem
    \item \textbf{Odpowiedz}: Model generuje tekst i zwraca uzytkownikowi
\end{enumerate}

\subsection{Diagram Przepywu}

\begin{verbatim}
Bankier.pl
    |
    | HTML
    v
Scraper (scraper.py)
    |
    | Symbole, Ceny
    v
DB Updater (db_updater.py)
    |
    | SQL INSERT/UPDATE
    v
MySQL Database
    |
    | SQL SELECT
    v
FastAPI (main.py)
    |
    | Dane + Prompt
    v
Ollama (Model AI)
    |
    | Tekst odpowiedzi
    v
Uzytkownik/Client
\end{verbatim}

\section{Interfejs API}

\subsection{Endpoint POST /ask}

Glowny endpoint do komunikacji z systemem.

\textbf{Request:}
\begin{lstlisting}[language=json]
POST /ask HTTP/1.1
Content-Type: application/json

{
  "prompt": "Porowniaj PKO i PKNORLEN"
}
\end{lstlisting}

\textbf{Response:}
\begin{lstlisting}[language=json]
{
  "response": "PKO Bank Polski handluje na poziomie 35.50 PLN, 
              podczas gdy PKNORLEN na 34.20 PLN. 
              PKO jest drozsze o 1.30 PLN...",
  "data": {...}
}
\end{lstlisting}

\subsection{Przyklady Zapytan}

\subsubsection{Porownanie Spółek}

\texttt{Zapytanie}: "Porowniaj PKO z PZU"

\texttt{Odpowiedz}: System pobiera ceny obu spółek z bazy i generuje porownanie oparte na danych.

\subsubsection{Przeglad Rynku}

\texttt{Zapytanie}: "Jaki jest stan rynku WIG20?"

\texttt{Odpowiedz}: Pobieranie wszystkich aktualnych cen i generowanie ogolnego przeglądu.

\subsubsection{Analiza Pojedynczej Spółki}

\texttt{Zapytanie}: "Analiza ALLEGRO"

\texttt{Odpowiedz}: Dane o spółce z historichnym kontekstem i analiza.

\section{Baza Danych}

\subsection{Schemat Tabel}

\subsubsection{Tabela Symbols}

\begin{table}[H]
\centering
\begin{tabular}{|l|l|l|}
\hline
\textbf{Kolumna} & \textbf{Typ} & \textbf{Opis} \\
\hline
id & INT PRIMARY KEY & Identyfikator unikatowy \\
name & VARCHAR(255) UNIQUE & Symbol spółki (PKO, PKNORLEN itp.) \\
market\_cap & FLOAT & Kapitalizacja rynkowa \\
\hline
\end{tabular}
\caption{Struktura tabeli Symbols}
\end{table}

\subsubsection{Tabela Prices}

\begin{table}[H]
\centering
\begin{tabular}{|l|l|l|}
\hline
\textbf{Kolumna} & \textbf{Typ} & \textbf{Opis} \\
\hline
id & INT PRIMARY KEY & Identyfikator unikatowy \\
symbol\_id & INT FOREIGN KEY & Odniesienie do Symbols \\
value & FLOAT & Cena akcji w PLN \\
time & DATETIME & Czas pobrania danych \\
\hline
\end{tabular}
\caption{Struktura tabeli Prices}
\end{table}

\subsection{Operacje na Bazie}

System wykorzystuje operacje:
\begin{itemize}
    \item \textbf{INSERT ... ON DUPLICATE KEY UPDATE} - idempotentne zapisywanie
    \item \textbf{UNIQUE(symbol\_id, time)} - unikanie duplikatów cen
    \item \textbf{FOREIGN KEY} - konsystencja danych
    \item \textbf{INDEX} - szybkie wyszukiwanie
\end{itemize}

\section{Bezpieczenstwo}

\subsection{Implementowane Mechanizmy}

\begin{itemize}
    \item \textbf{Hasła w .env}: Dane wrażliwe nie sa w kodzie
    \item \textbf{Health Checks}: Sprawdzenie gotowosci serwisów
    \item \textbf{Timeout}: Limity czasowe na polaczenia (10-20 sekund)
    \item \textbf{Error Handling}: Logowanie bledów z kontekstem
    \item \textbf{UNIQUE Constraints}: Zapobieganie duplikatom
    \item \textbf{User-Agent Emulation}: Unikanie blokady przez scraper
\end{itemize}

\subsection{Zalecenia dla Produkcji}

\begin{itemize}
    \item Wdrozenie autentykacji API (JWT lub API keys)
    \item Dodanie rate limiting na endpoint /ask
    \item Szyfrowanie polaczen do MySQL (SSL/TLS)
    \item Regularne aktualizacje zaleznosci Python
    \item Monitoring i logowanie akcji uzytkownikow
    \item Regularne backupy bazy danych
    \item Wdrozenie WAF (Web Application Firewall)
\end{itemize}

\section{Instalacja i Urzadzanie}

\subsection{Kroki Instalacji}

\begin{enumerate}
    \item Klonowanie repozytorium z GitHub
    \item Konfiguracja pliku .env z hasłami
    \item Budowa kontenerow Docker
    \item Uruchomienie Docker Compose
    \item Pobranie modelu Ollama (qwen2.5)
    \item Weryfikacja dostepnosci API na http://localhost:8000
\end{enumerate}

\subsection{Komendy}

\begin{lstlisting}[language=bash]
# Uruchomienie
docker-compose up -d

# Zatrzymanie
docker-compose down

# Logi API
docker logs wig20_api -f

# Logi MySQL
docker logs wig20_db -f

# Status kontenow
docker-compose ps
\end{lstlisting}

\subsection{Dostep do Serwisów}

\begin{itemize}
    \item \textbf{FastAPI}: http://localhost:8000
    \item \textbf{Swagger UI}: http://localhost:8000/docs
    \item \textbf{ReDoc}: http://localhost:8000/redoc
    \item \textbf{Ollama}: http://localhost:11434
    \item \textbf{MySQL}: localhost:3306 (uzytkownik: agent)
\end{itemize}

\section{Monitorowanie i Utrzymanie}

\subsection{Logowanie}

Wszystkie komponenty loguja do stdout / stderr:

\begin{itemize}
    \item API loguje zapytania, bledy, czasy odpowiedzi
    \item Scraper loguje status pobrania i bledy polaczenia
    \item DB Updater loguje operacje SQL
    \item Ollama loguje inference requests
\end{itemize}

\subsection{Czyszczenie Bazy Danych}

Zalecane operacje:

\begin{lstlisting}[language=sql]
-- Czyszczenie danych starszych niz 30 dni
DELETE FROM Prices 
WHERE time < DATE_SUB(NOW(), INTERVAL 30 DAY);

-- Optymalizacja tabeli
OPTIMIZE TABLE Prices;
OPTIMIZE TABLE Symbols;

-- Sprawdzenie rozmiaru
SHOW TABLE STATUS LIKE 'Prices';
\end{lstlisting}

\section{Przyszłe Rozszerzenia}

\subsection{Krotkoterminowe (1-2 miesiace)}

\begin{itemize}
    \item Alerty cenowe (email/SMS/push notification)
    \item Cache Redis dla czestych zapytan
    \item Filtry czasowe w /ask ("cena z ostatniego miesiaca")
    \item Wsparcie dla innych indeksów (WIG, mWIG40, sWIG80)
    \item Statystyki i wykresy cenowe
\end{itemize}

\subsection{Dlugotermionowe (3-6 miesiecy)}

\begin{itemize}
    \item Machine Learning dla predykcji cen
    \item Frontend webowy (React/Vue.js)
    \item Aplikacja mobilna (iOS/Android)
    \item Integracja z platformami tradingowymi
    \item Raportowanie i eksport danych (CSV, PDF)
    \item Dashboard z wizualizacją danych
    \item Wsparcie dla portfeli uzytkowników
\end{itemize}

\section{Struktura Projektu}

\begin{verbatim}
wig20-agent/
├── main.py              # FastAPI server
├── scraper.py           # Web scraping
├── db_updater.py        # Database operations
├── mcp_server.py        # MCP protocol server
├── client.py            # Client implementation
├── requirements.txt     # Python dependencies
├── Dockerfile           # Container image
├── docker-compose.yml   # Multi-container setup
├── REPORT.md            # Markdown documentation
├── RAPORT.tex           # LaTeX report
├── .env.example         # Environment template
├── .gitignore          # Git ignore rules
└── __pycache__/        # Python cache
\end{verbatim}

\section{Podsumowanie}

WIG20 Agent to kompleksowy, skalowany system do inteligentnej analizy danych gieldowych. Architektura oparta na mikroserwisach i konteneryzacji zapewnia łatwosc wdrozenia i utrzymania. Integracja z modelami AI umozliwia interaktywne odpowiadanie na pytania uzytkownikow w jezyku naturalnym.

System automatycznie:
\begin{itemize}
    \item Pobiera i aktualizuje dane o cenach akcji
    \item Przechowuje dane w strukturyzowanej bazie
    \item Odpowiada na pytania uzytkownikow
    \item DostarczaAPI REST do integracji
\end{itemize}

Projekt ma potencjal do dalszego rozwoju w kierunku zaawansowanej analizy finansowej, predykcji trendow i wsparcia dla decyzji inwestycyjnych.

\end{document}

\section{Funkcjonalność}

System oferuje:
\begin{itemize}
    \item \textbf{Pobieranie}: Automatyczne pobranie danych z Bankier.pl (symbole, ceny, kapitalizacja)
    \item \textbf{Baza danych}: MySQL przechowuje spółki i ceny historyczne
    \item \textbf{API}: Endpoint POST /ask do wysyłania pytań
    \item \textbf{AI}: Ollama odpowiada na pytania w języku polskim i angielskim
\end{itemize}

\section{Architektura Systemu}

\subsection{Diagram architektury}

\begin{verbatim}
    ┌─────────────┐
    │ Web Scraper │
    │ (scraper.py)│
    └──────┬──────┘
           │
           │ Dane o cenach
           ▼
    ┌─────────────────┐
    │ Database Updater│
    │(db_updater.py)  │
    └──────┬──────────┘
           │
           ▼
    ┌─────────────────┐
    │  MySQL 8.0      │
    │  Database       │
    └────┬────────┬───┘
         │        │
         │        │
    ┌────▼───┐  ┌─▼─────────┐
    │FastAPI │  │MCP Server │
    │API     │  │(mcp_server)
    │main.py │  └─┬──────────┘
    └────┬───┘    │
         │        │
         │ Zapytania
         │        │
    ┌────▼────────▼──}

Komponenty systemu:
\begin{enumerate}
    \item \textbf{scraper.py} - pobiera dane z Bankier.pl (BeautifulSoup)
    \item \textbf{db\_updater.py} - zapisuje dane do MySQL
    \item \textbf{main.py} - REST API (FastAPI) na porcie 8000
    \item \textbf{mcp\_server.py} - Model Context Protocol Server
\end{enumerate}

Przepływ: Web Scraper → MySQL ← FastAPI + Ollama ← Clientxtbf{Docker Compose}: 1.29+
    \item \textbf{RAM}: 4GB minimum
    \item \textbf{Dysk}: 10GB dla bazy danych i modelu Ollama
\end{itemize}

\subsection{Zależności Python}

\begin{lstlisting}[caption=requirements.txt]
fastapi          # Web framework
uvicorn          # ASGI server
mysql-connector-python  # MySQL driver
python-dotenv    # Environment variables
requests         # HTTP client
pydantic         # Data validation
beautifulsoup4   # HTML parsing
mcp              # Model Context Protocol
starlette        # ASGI toolkit
\end{lstlisting}

\section{Architektura Deployment}

\subsection{Docker Compose}

System uruchamia się w czterech kontenerach Docker:

\begin{table}[H]
\centering
\begin{tabular}{|l|l|l|l|}
\hline
\textbf{Serwis} & \textbf{Obraz} & \textbf{Port} & \textbf{Opis} \\
\hline
db & mysql:8.0 & 3306 & Baza danych MySQL \\
ollama & ollama/ollama & 11434 & Model AI \\
api & local build & 8000 & FastAPI service \\
mcp & local build & - & MCP Server \\
\hline
\end{tabular}
\caption{Serwisy Docker Compose}
\end{table}

\subsection{Health Check i Dependency Ordering}

\begin{itemize}
    \item \textbf{DB Health Check}: \texttt{mysqladmin ping}
    \item \textbf{Timeout}: 20 sekund
    \item \textbf{Retries}: 10 prób
    \item \textbf{Zależności}:
    \begin{itemize}
        \item API czeka na gotowość DB
        \item API czeka na start Ollama
        \item MCP czeka na gotowość DB
    \end{itemize}
\end{itemize}

\subsection{Zmienne środowiskowe}

Konfiguracja pobierana z pliku \texttt{.env}:

\begin{lstlisting}[caption=Zmienne .env]
MYSQL_ROOT_PASSWORD=xxx
MYSQL_DATABASE=wig20
MYSQL_USER=agent
MYSQL_PASSWORD=xxx
\end{lstlisting}

\section{Przepływ Danych}

\begin{enumerate}
    \item Scraper pobiera dane z Bankier.pl
    \item db\_updater zapisuje do MySQL
    \item Użytkownik wysyła POST /ask
    \item API pobiera dane z bazy SQL
    \item Ollama generuje odpowiedź
    \item Odpowiedź wysłana do użytkownika
\end{enumerate}

\section{Interfejs Użytkownika}

\subsection{Przykładowe zapytania}

\subsubsection{Porównanie spółek}

\begin{lstlisting}[caption=Przykład 1: Porównanie]
Request:
POST /ask
{"prompt": "Porównaj PKO i PKNORLEN"}

Response:
"PKO Bank Polski: 35.50 PLN
PKOBP (PKO Banku Polskiego): 34.20 PLN
PKO jest droższe od PKNORLEN o około 1.30 PLN..."
\end{lstlisting}

\subsubsePrzykładowe Zastosowania}

\textbf{Porównanie}: POST /ask z \texttt{"prompt": "Porównaj PKO i PKNORLEN"}

\textbf{Przegląd}: POST /ask z \texttt{"prompt": "Jaki jest stan rynku?"}

\textbf{Analiza}: POST /ask z \texttt{"prompt": "Analiza PKO"
    \item Wdrożyć autentykację API (JWT, API Key)
    \item Dodać rate limiting na endpoint /ask
    \item Szyfrować połączenia do MySQL (SSL/TLS)
    \item Regularnie aktualizować zależności Python
    \item Monitorować logi w produkcji
    \item Backupować bazę danych MySQL
\end{itemize}

\section{Testowanie}

\subsection{Scenariusze testowe}

\subsubsection{Test jednostkowy: Ekstrakcja symboli}

\begin{lstlisting}[caption=Test ekstrakcji symboli]
extract_sTechnologie}

\begin{itemize}
    \item Python 3.11 - język programowania
    \item FastAPI + Uvicorn - REST API
    \item MySQL 8.0 - baza danych
    \item BeautifulSoup4 - web scraping
    \item Ollama + Qwen2.5 - model AI
    \item Docker Compose - konteneryzacja
\end{itemize}

Wymagania: Python 3.11+, Docker, 4GB RAM.env
# Edytować .env z właściwymi hasłami
    \end{lstlisting}

    \item \textbf{Budowa i uruchomienie}:
    \begin{lstlisting}[language=bash]
docker-coDeployment}

Docker Compose uruchamia 4 serwisy:
\begin{itemize}
    \item db (MySQL 8.0)
    \item ollama (Model AI)
    \item api (FastAPI na porcie 8000)
 begin{itemize}
    \item Hasła w zmiennych .env, nie w kodzie
    \item Health checks dla serwisów
    \item Error handling z logowaniem
    \item UNIQUE constraints w bazie
\end{itemize}

Zalecenia: Autentykacja API (JWT), rate limiting, SSL/TLS, regularne aktualizacje, backupye}
    \item Dodanie alertów cenowych (integracja email/SMS)
    \item Implementacja cache'u Redis
    \item Dodanie filtrów w /ask (data, zakres)
    \item Wsparcie dla więcej indeksów (WIG, mWIG40, sWIG80)
\end{itemize}

\subsection{Długoterminowe}

\begin{itRozszerzenia}

Krótkoterminowe: Alerty cenowe, cache Redis, filtry w /ask, więcej indeksów

Długoterminowe: Predykcja ML, frontend webowy, aplikacja mobilna, dashboardmpose.yml   # Service orchestration
├── REPORT.md            # Markdown report
├── RAPORT.tex          # LaTeX report (ten dokument)
├── __pycache__/        # Python cache
└── .env                # Configuration (not in repo)
\end{verbatim}

\section{Zaključenie}

Projekt WIG20 Agent stanowi nowoczesny i funkcjonalny system do zarządzania danymi giełdowymi. Architektura mikroserwisów z Docker Compose zapewnia skalowlaność i łatwość wdrażania. Integracja z AI (Ollama) umożliwia użytkownikom interaktywne korzystanie z danych poprzez pytania w języku naturalnym.

System jest w stanie:
\begin{itemize}
    \item Automatycznie pobierać i aktualizować dane giełdowe
    \item Przechowywać dane historyczne dla analiz
    \item Udostępniać API RESTful do integracji
    \item Generować inteligentne odpowiedzi na zapytania użytkowników
\end{itemize}
}

\begin{verbatim}
main.py, scraper.py, db_updater.py, mcp_server.py, client.py
requirements.txt, Dockerfile, docker-compose.yml, .env
# RestartPodsumowanie}

WIG20 Agent to system do pobierania i analizy danych giełdowych z indeksu WIG20. Architektura mikroserwisów zapewnia skalowlaność, a integracja z AI umożliwia interaktywne odpowiadanie na pytania w języku naturalnym. System automatycznie pobiera dane, przechowuje je w bazie i udostępnia poprzez REST API.